% BHU PYQ -- Auto-generated & error-corrected
\documentclass[11pt,a4paper]{article}

\usepackage[a4paper,top=2.5cm,bottom=2.5cm,left=2.8cm,right=2.8cm,headheight=14pt]{geometry}
\usepackage{amsmath,amssymb}
\usepackage[T1]{fontenc}
\usepackage[utf8]{inputenc}
\usepackage{lmodern}
\usepackage{microtype}
\usepackage{fancyhdr}
\usepackage{enumitem}
\usepackage{hyperref}
\usepackage{lastpage}
\usepackage{parskip}

\hypersetup{colorlinks=true,urlcolor=black,linkcolor=black,
  pdftitle={BSC-07A: Descriptive Statistics -- BHU PYQ},pdfauthor={Banaras Hindu University}}

\pagestyle{fancy}\fancyhf{}
\renewcommand{\headrulewidth}{0.4pt}\renewcommand{\footrulewidth}{0.4pt}
\fancyhead[L]{\small   \textit{Banaras Hindu University}}
\fancyhead[R]{\small   \textit{BSC-07A: Descriptive Statistics}}
\fancyfoot[C]{\small Page  \thepage\ of \pageref{LastPage}}


\newlist{parts}{enumerate}{2}
\setlist[parts,1]{label=(\alph*),leftmargin=*,itemsep=5pt,topsep=3pt}
\setlist[parts,2]{label=(\roman*),leftmargin=*,itemsep=3pt,topsep=2pt}

\newcommand{\pts}[1]{\hfill   \textit{[#1]}}


\begin{document}


\begin{center}
  {\large\bfseries BANARAS HINDU UNIVERSITY}\\[2pt]
  {
\normalsize B.Sc. (Hons.) Semester II Examination, 2013-14}\\[6pt]
  \rule{0.8\linewidth}{0.5pt}\\[6pt]
  {\large\bfseries Paper No. BSC-07A: Descriptive Statistics}\\[4pt]
  {
\normalsize Time: 3 Hours\hfill Maximum Marks: 70}
\end{center}
\rule{\linewidth}{0.4pt}
\smallskip



\noindent   \textit{Instructions: Answer   \textbf{five} questions including Question~1 which is compulsory. Symbols have their usual meanings.}
\bigskip

\medskip


\noindent   \textbf{Question 1.}

\begin{parts}
  \item Write down the definition of Statistics. Explain in brief the nature and scope of statistics along with its limitations. \pts{7}
  \item Explain with examples nominal scale, ordinal scale, interval scale, and ratio scale of measurement. \pts{7}
\end{parts}
\medskip\rule{\linewidth}{0.2pt}

\medskip


\noindent   \textbf{Question 2.}

\begin{parts}
  \item What do you mean by primary and secondary data? Write in brief the sources of secondary data. Describe the schedule method of collecting primary data. \pts{7}
  \item Describe various diagrammatical methods for the representation of data. \pts{7}
\end{parts}
\medskip\rule{\linewidth}{0.2pt}

\medskip


\noindent   \textbf{Question 3.}
Describe histogram, frequency polygon, frequency curve, and ogive. Draw the histogram and cumulative frequency curve (ogive) for the following data and find the median value with the help of the ogive:

\begin{center}

\begin{tabular}{|l|c|c|c|c|c|c|}
\hline
   \textbf{Class Interval} & 5--7 & 7--9 & 9--11 & 11--13 & 13--15 & 15--17 \\ \hline
   \textbf{Frequency}      & 6    & 50   & 80    & 55     & 20     & 15     \\ \hline
\end{tabular}
\end{center}
\pts{14}
\medskip\rule{\linewidth}{0.2pt}

\medskip


\noindent   \textbf{Question 4.}
Define arithmetic mean, median, and mode of a frequency distribution and discuss their merits and demerits. Discuss the effect of change of origin and scale on arithmetic mean. \pts{14}
\medskip\rule{\linewidth}{0.2pt}

\medskip


\noindent   \textbf{Question 5.}
What do you mean by dispersion? Explain various measures of dispersion with their relative merits and demerits. Define the coefficient of variation and explain why it is required. \pts{14}
\medskip\rule{\linewidth}{0.2pt}

\medskip


\noindent   \textbf{Question 6.}

\begin{parts}
  \item Define Quartile, Decile, and Percentile of a frequency distribution. Calculate $Q_1$, $D_5$, and $P_{90}$ for the following distribution:
  
\begin{center}
  
\begin{tabular}{|l|c|c|c|c|c|}
  \hline
   \textbf{Class}     & 20--30 & 30--40 & 40--50 & 50--60 & 60--70 \\ \hline
   \textbf{Frequency} & 12     & 25     & 40     & 15     & 8      \\ \hline
  \end{tabular}
  \end{center} \pts{7}
  \item Define raw and central moments for grouped data. Explain Skewness and Kurtosis and discuss their various measures. \pts{7}
\end{parts}
\medskip\rule{\linewidth}{0.2pt}

\medskip


\noindent   \textbf{Question 7.}
What is a scatter diagram? Define the correlation coefficient and give examples where positive and negative correlation coefficients exist. Show that the value of the correlation coefficient lies between $-1$ and $+1$. \pts{14}
\medskip\rule{\linewidth}{0.2pt}

\medskip


\noindent   \textbf{Question 8.}

\begin{parts}
  \item What do you mean by regression? Write the expression for the two regression lines. Given the two regression equations:
  \[
  8X - 10Y = -46.6 \quad  \text{and} \quad 40X - 18Y = 214
  \]
  calculate the mean values of $X$ and $Y$, and the correlation coefficient between them. \pts{7}
  \item Ten competitors in a musical contest were ranked by three judges $A$, $B$, and $C$ in the following order:
  
\begin{center}
  
\begin{tabular}{|l|cccccccccc|}
  \hline
   \textbf{Judge A} & 1 & 6 & 5 & 10 & 3 & 2 & 4 & 9 & 7 & 8 \\
   \textbf{Judge B} & 3 & 5 & 8 & 4  & 7 & 10& 2 & 1 & 6 & 9 \\
   \textbf{Judge C} & 6 & 4 & 9 & 8  & 1 & 2 & 3 & 10& 5 & 7 \\ \hline
  \end{tabular}
  \end{center}
  Using the rank correlation method, discuss which pair of judges has the nearest approach to common likings in music. \pts{7}
\end{parts}

\vfill
\rule{\linewidth}{0.4pt}

\begin{center}\small
  \href{https://bhu-exam.vercel.app/}{Click here to visit our website}\\[4pt]
  \href{https://me-aryan.vercel.app/}{Click here to visit our contributor}\\[4pt]
  \href{https://quashan-me.vercel.app/}{Click here to visit our contributor}
\end{center}

\end{document}